\section{Introduction}

The standard agent loop pays for a model decision at every tool
boundary~\citep{yao2023react}. That flexibility is valuable while a workflow is
novel, but mature agents often revisit the same read-only evidence path with the
same data dependencies, where repeating the decision adds latency, tokens, and
cost without adding adaptation.

Removing that decision is not ordinary caching. A tool argument may depend on a
specific earlier observation, a repeated sequence may cross an approval or
write, an apparently read-only tool may hide an undeclared effect, and replaying
the same outputs may still change the downstream answer. The problem is not
``find recurring traces'' but deciding when a recurrent model-mediated region
may be replaced by a deterministic program, and when it must be refused.

Consider held-out issue \#4420. Given only \code{issue\_number=4420}, the
observed trace calls \code{record(4420)}, passes the \emph{returned}
\code{record.issue\_number} to \code{labels} and \code{comments}, and invokes
the latter with \code{limit=100}. All 132 discovery executions take that route,
so a recurrence-only optimizer would emit all three reads. \method instead
retires the \emph{region}: \code{comments.limit} is reachable neither from the
entry state nor from either earlier result, so one unwitnessed slot blocks the
three-read candidate as \code{ungroundable\_slot}. What survives is the maximal
justified prefix \code{record} $\rightarrow$ \code{labels}; the agent still
chooses the limit and writes the answer. That two-read artifact passes 92/92
calibration groups at $\alpha=0.05$ (upper bound 0.0498).

An earlier 18-record study supplies the counterexample: its
three-read artifact replayed all 45 tool calls, yet issue \#6602 lost the URL
from a Markdown link and only 17/18 downstream answers stayed exact. A
provider-free continuation guard detects and checked-renders that miss --- a
mechanism check, not live safety evidence. Together the two cohorts
(\cref{fig:aha-example}) show why recurrence only proposes a candidate.

\begin{figure}[t]
  \centering
  \includegraphics[width=0.72\linewidth]{figures/gac_aha_example.pdf}
  \Description{A real GitHub issue trace repeats record, labels, and comments, but GAC
  rejects the full candidate because the comments limit is ungroundable and emits only
  the record and labels prefix.  A separate issue 6602 continuation miss shows why clean
  tool replay does not certify the final answer.}
  \caption{Two retained cohorts expose different failure modes: issue \#4420
forces a shorter grounded prefix, and issue \#6602 shows that clean tool replay
need not preserve the downstream answer.}
  \label{fig:aha-example}
\end{figure}

The closest precedent is the tracing just-in-time compiler:
Dynamo~\citep{bala2000dynamo} detects a hot path, compiles it behind a
\emph{guard}, and deoptimizes when the guard fails~\citep{holzle1992deopt}.
Agent traces admit the same structure --- hot-trace detection becomes family
mining, guard synthesis an entry contract, deoptimization fallback to the
unchanged agent --- but add three obligations, and claim one capability less: a
JIT enters a trace at an arbitrary program point, whereas \method compiles only
\emph{initial} read-only prefixes (\cref{eq:prefix}), because its runtime
resolves artifacts at the entry boundary. The obligations are that tool
arguments be shown \emph{grounded} in observable state rather than invented;
that declared \emph{effects} permit speculation, since an agent's
``interpreter'' can call the outside world; and that the residual error rate,
being statistical, needs a finite-sample bound rather than a soundness proof.

We study these obligations through \emph{guarded agentic compaction} (\method),
a compile-or-retire trace-to-program system for recurrent read-only prefixes
(\cref{fig:pipeline}). The core contribution is not ``agent compilation'' in
general, but an admissibility argument for trace-derived specialization:

\begin{enumerate}[leftmargin=1.2em,itemsep=1pt,topsep=2pt]
  \item A guarded specialization pipeline (\cref{alg:compile}) combining typed
provenance, effect and position barriers, bounded synthesis, runtime
verification, and finite-sample compile-or-retire admission
(\cref{alg:calibrate}, \cref{prop:admission}) with the certificate level stated
per artifact; every stage can refuse, and the default output is the unchanged
agent. \item A real-provider evaluation across
three public-record workflow families, plus a narrower cross-repository, time-forward extension keeping both artifacts and principled retirements.
  \item A negative-result boundary: recurrence and replayability do not license
  deployment. NESTFUL, API-Bank, and executed BFCL gold plans retire under the
  same admission logic, and hand-written programs are competitive or tied wherever
  the workflow is already known.
  \item A pre-registered fourth substrate, AppWorld --- which admits --- plus a
  measurement on 8{,}190 released baseline trajectories where the admitted region
  is eligible on 2{,}339/2{,}340 full-code but only 1/4{,}680 ReAct or
  plan-and-execute trajectories, separating \emph{admissible} from
  \emph{dispatchable}.
\end{enumerate}
